\documentclass[aspectratio=1610, 9pt]{beamer}
\usepackage{mobi-beamer}

\title[Congruencia lineal y transformada i…]{Congruencia lineal y transformada inversa}
\subtitle{Alejandra Tabares}
\author{Alejandra Tabares Pozos}
\institute{Universidad de los Andes \\
 Departamento de Ingeniería Industrial}
\date{}

\begin{document}

% -----------------------------------------------------------------------------
% SECCIÓN 1: EL GANCHO
% -----------------------------------------------------------------------------
\begin{frame}
    \titlepage
\end{frame}

\begin{frame}{El Caso de Estudio: Auditoría al Casino Andino}
    \begin{columns}
        \column{0.6\textwidth}
        \textbf{El problema}\\
        Una plataforma de apuestas en línea, ``Casino Andino'', ha sido acusada de fraude.
        Jugadores afirman que las tragamonedas son predecibles y que los premios grandes
        no salen según las probabilidades anunciadas.

        \vspace{0.45cm}

        \textbf{Tu misión}\\
        Como ingenieras e ingenieros, debemos auditar el \textbf{motor de simulación}:
        \begin{itemize}
            \item ¿Cómo genera el computador la ``suerte''?
            \item ¿Qué significa ``aleatorio'' en un sistema determinístico?
            \item ¿Cómo se traducen números uniformes en premios y tiempos de juego?
        \end{itemize}

        \column{0.4\textwidth}
        \centering
        \textbf{¿Azar o algoritmo?}
        \begin{tikzpicture}
            \node[draw, circle, minimum size=2cm, fill=red!20] (A) at (0,2) {7};
            \node[draw, circle, minimum size=2cm, fill=blue!20] (B) at (0,0) {Bar};
            \node[draw, circle, minimum size=2cm, fill=yellow!20] (C) at (0,-2) {Cereza};
            \draw[-{Stealth[length=3mm]}, thick] (-2, 0) -- node[above] {\small Entrada} (-1,0);
        \end{tikzpicture}
    \end{columns}
\end{frame}

\begin{frame}{La Paradoja del Computador Determinístico}
    \begin{alertblock}{Obstáculo fundamental}
    Un computador es una máquina de estados finitos y determinísticos.
    Si le pides que calcule $2+2$, siempre dará $4$.
    \end{alertblock}

    \vspace{0.45cm}

    \textbf{Pregunta}\\
    Si el hardware es determinístico, ¿cómo podemos simular incertidumbre estocástica
    para modelar el juego?

    \vspace{0.25cm}
    \textit{Necesitamos un mecanismo matemático que produzca secuencias con comportamiento estadístico parecido al azar.}
\end{frame}

% -----------------------------------------------------------------------------
% SECCIÓN 2: ANATOMÍA DE LA SIMULACIÓN
% -----------------------------------------------------------------------------
\section{Anatomía de la simulación}

\begin{frame}{Definición Formal de Simulación Estocástica}
    Antes de auditar el código, fijemos el marco. Según \textbf{Banks et al.}:

    \begin{quote}
        La simulación es la imitación de la operación de un proceso o sistema a través del tiempo,
        generando una historia artificial para inferir características operacionales.
    \end{quote}

    \vspace{0.25cm}
    \textbf{Diferencia clave}
    \begin{itemize}
        \item \textbf{Optimización determinística:} busca $x^\star$ analíticamente.
        \item \textbf{Simulación estocástica:} estima cantidades como $E[Y]$ a partir de repeticiones.
    \end{itemize}
\end{frame}

\begin{frame}{Los Elementos Constitutivos del Modelo}
    Para el casino, formalizamos:
    \[
    Y=f(X_1,X_2,\dots,X_n).
    \]

    \vspace{0.2cm}
    \begin{enumerate}
        \item \textbf{Entradas $X_i$}\\
        Variables aleatorias que impulsan el sistema.
        \begin{itemize}
            \item Ejemplos: tiempo entre apuestas, monto apostado, evento de premio grande.
        \end{itemize}

        \item \textbf{Mecanismo $f$}\\
        La lógica del proceso.
        \begin{itemize}
            \item Ejemplos: reglas del tragamonedas, tabla de pagos, validaciones del servidor.
        \end{itemize}

        \item \textbf{Salidas $Y$}\\
        Métricas de desempeño que también son aleatorias.
        \begin{itemize}
            \item Ejemplos: ganancia neta diaria, retorno al jugador, frecuencia de premios grandes.
        \end{itemize}
    \end{enumerate}
\end{frame}

\begin{frame}{El Algoritmo General de Simulación}
    \begin{center}
    \begin{tikzpicture}[node distance=1.35cm]
        \node (start) [draw, rectangle, rounded corners] {Inicio};
        \node (rng) [draw, rectangle, below=of start, fill=blue!10] {1. Generar $U \in [0,1)$};
        \node (variate) [draw, rectangle, below=of rng, fill=yellow!20] {2. Transformar $U$ en $X$};
        \node (model) [draw, rectangle, below=of variate] {3. Evaluar $Y=f(X)$};
        \node (stats) [draw, rectangle, below=of model] {4. Recolectar estadísticas};
        \node (end) [draw, rectangle, rounded corners, below=of stats] {Análisis de salida};

        \draw[-{Stealth[length=3mm]}] (start) -- (rng);
        \draw[-{Stealth[length=3mm]}] (rng) -- (variate);
        \draw[-{Stealth[length=3mm]}] (variate) -- (model);
        \draw[-{Stealth[length=3mm]}] (model) -- (stats);
        \draw[-{Stealth[length=3mm]}] (stats) -- (end);

        \node [right=2.6cm of rng, text width=4.8cm] {\small \textbf{Hoy}\\De dónde sale $U$\\Aritmética de reloj y LCG};
        \node [right=2.6cm of variate, text width=4.8cm] {\small \textbf{Hoy}\\Cómo pasar de $U$ a $X$\\Transformada inversa};
    \end{tikzpicture}
    \end{center}
\end{frame}

% -----------------------------------------------------------------------------
% SECCIÓN 3: ARITMÉTICA DE RELOJ
% -----------------------------------------------------------------------------
\section{Aritmética de reloj}

\begin{frame}{Imagen mental: aritmética de reloj}
    \textbf{Pensemos en un reloj de 12 horas}

    \vspace{0.35cm}
    \begin{itemize}
        \item Si son las 10 y pasan 5 horas, no dices 15; dices 3.
        \item La razón es simple: al llegar a 12, el reloj vuelve a empezar.
        \item El reloj no guarda el número total, guarda la \textbf{posición} después de dar vueltas completas.
    \end{itemize}

    \vspace{0.35cm}
    \begin{block}{Idea que vamos a formalizar}
    ``Quedarse con la posición'' equivale a \textbf{quedarse con el residuo} al dividir por el tamaño del reloj.
    \end{block}
\end{frame}

\begin{frame}{Residuo y operador módulo}
    Cuando dividimos un entero $a$ entre un entero positivo $m$, existe un único par de enteros $(q,r)$ tal que:
    \[
    a = qm + r,
    \qquad 0\le r \le m-1.
    \]

    \vspace{0.25cm}
    \begin{block}{Definición}
    Llamamos \textbf{residuo} a $r$ y lo escribimos como:
    \[
    r = a \bmod m.
    \]
    \end{block}

    \vspace{0.25cm}
    \begin{exampleblock}{Ejemplo con $m=12$}
        \[
        15 = 1\cdot 12 + 3 \Rightarrow 15\bmod 12 = 3,
        \qquad
        27 = 2\cdot 12 + 3 \Rightarrow 27\bmod 12 = 3.
        \]
    \end{exampleblock}
\end{frame}

\begin{frame}{Congruencia como ``mismo residuo''}
    Diremos que dos números son \textbf{iguales en el reloj de tamaño $m$} si dejan el mismo residuo al dividir por $m$.

    \vspace{0.25cm}
    \begin{block}{Definición equivalente}
    \[
    a \equiv b \pmod m
    \quad \Longleftrightarrow \quad
    a \bmod m = b \bmod m.
    \]
    \end{block}

    \vspace{0.25cm}
    \textbf{Regla práctica}
    \[
    a \equiv b \pmod m
    \quad \Longleftrightarrow \quad
    a-b \text{ es múltiplo de } m.
    \]

    \vspace{0.15cm}
    \begin{exampleblock}{Ejemplo con $m=12$}
        $15\equiv 3\pmod{12}$ porque $15\bmod 12=3\bmod 12=3$ (y $15-3=12$).
    \end{exampleblock}
\end{frame}

\begin{frame}{Operaciones en el reloj}
    En aritmética de reloj operamos en enteros, pero el resultado final se reduce con $\bmod m$.

    \vspace{0.25cm}
    \begin{block}{Suma y multiplicación}
    \[
    (a+b)\bmod m \quad \text{y} \quad (ab)\bmod m
    \]
    son las posiciones finales del reloj después de sumar o multiplicar.
    \end{block}

    \vspace{0.25cm}
    \begin{exampleblock}{Ejemplo de suma en reloj 12}
        \[
        (10+5)\bmod 12 = 15\bmod 12 = 3.
        \]
    \end{exampleblock}

    \vspace{0.15cm}
    \begin{exampleblock}{Ejemplo de multiplicación en reloj 7}
        \[
        (6\cdot 5)\bmod 7 = 30\bmod 7 = 2.
        \]
    \end{exampleblock}
\end{frame}

\begin{frame}{El reloj produce estados en $\{0,1,\dots,m-1\}$}
    En simulación vamos a trabajar con un \textbf{estado} $Z$ que siempre viva dentro del reloj.

    \vspace{0.25cm}
    \begin{block}{Hecho fundamental}
    Para cualquier entero $t$,
    \[
    Z = t \bmod m \quad \Rightarrow \quad Z\in\{0,1,\dots,m-1\}.
    \]
    \end{block}

    \vspace{0.25cm}
    \begin{exampleblock}{Ejemplo con $m=12$}
        \[
        0\bmod 12=0,\ 1\bmod 12=1,\ \dots,\ 11\bmod 12=11,\ 12\bmod 12=0,\ 13\bmod 12=1.
        \]
        El reloj obliga a ``volver'' cuando se completa una vuelta.
    \end{exampleblock}
\end{frame}

\begin{frame}{Una secuencia de estados usando el reloj}
    Un solo estado no basta: necesitamos una \textbf{secuencia} $Z_0,Z_1,Z_2,\dots$.

    \vspace{0.25cm}
    \begin{block}{Regla mínima para entender el mecanismo}
    \[
    Z_{n+1} = (Z_n+1)\bmod m.
    \]
    \end{block}

    \vspace{0.2cm}
    \begin{exampleblock}{Ejemplo con $m=10$ y $Z_0=7$}
        \[
        7,\ 8,\ 9,\ 0,\ 1,\ 2,\ 3,\ 4,\ 5,\ 6,\ 7,\dots
        \]
    \end{exampleblock}

    \vspace{0.2cm}
    \begin{alertblock}{Lectura}
    Ya tenemos lo esencial: \textbf{actualizar} y luego \textbf{recortar con $\bmod m$}.
    Lo que falta es que la secuencia no sea un patrón obvio.
    \end{alertblock}
\end{frame}

\begin{frame}{De estados $Z_n$ a números $U_n$ en $[0,1)$}
    En simulación, el algoritmo general comienza generando $U\in[0,1)$.

    \vspace{0.25cm}
    \begin{block}{Normalización}
    Si $Z_n\in\{0,1,\dots,m-1\}$, definimos:
    \[
    U_n=\frac{Z_n}{m}.
    \]
    Entonces $U_n\in[0,1)$.
    \end{block}

    \vspace{0.2cm}
    \begin{exampleblock}{Ejemplo con $m=10$}
        \[
        Z_n: 7,\ 8,\ 9,\ 0,\ 1,\dots
        \quad\Rightarrow\quad
        U_n: 0.7,\ 0.8,\ 0.9,\ 0.0,\ 0.1,\dots
        \]
    \end{exampleblock}
\end{frame}

% -----------------------------------------------------------------------------
% SECCIÓN 4: ECUACIONES LINEALES EN EL RELOJ
% -----------------------------------------------------------------------------
\section{Ecuaciones lineales en el reloj}

\begin{frame}{La pregunta central}
    Ahora viene lo importante para auditoría y para entender generadores.

    \vspace{0.35cm}
    \begin{block}{Problema}
    Encontrar $x$ tal que
    \[
    ax \equiv b \pmod m.
    \]
    \end{block}

    \vspace{0.25cm}
    \textbf{Traducción en español}
    \begin{itemize}
        \item Multiplico $a\cdot x$.
        \item Me quedo con el residuo al dividir por $m$.
        \item Quiero que esa posición sea $b$.
    \end{itemize}
\end{frame}

\begin{frame}{Ejemplo guiado en reloj 12}
    Resolver:
    \[
    5x \equiv 1 \pmod{12}.
    \]

    \vspace{0.2cm}
    Interpretación: ``busco un $x$ para que $5x$ termine en la posición 1 del reloj 12.''

    \vspace{0.25cm}
    Probamos valores pequeños:
    \[
    \begin{array}{c|ccccc}
    x & 1 & 2 & 3 & 4 & 5\\ \hline
    5x & 5 & 10 & 15 & 20 & 25\\
    5x \bmod 12 & 5 & 10 & 3 & 8 & 1
    \end{array}
    \]

    \vspace{0.2cm}
    \begin{exampleblock}{Conclusión}
    $x=5$ funciona. En el reloj 12, multiplicar por 5 puede deshacerse multiplicando por 5 otra vez.
    \end{exampleblock}
\end{frame}

\begin{frame}{Inverso en el reloj}
    El ejemplo anterior introduce una idea clave:

    \vspace{0.25cm}
    \begin{block}{Definición}
    Decimos que $a$ tiene inverso en el reloj $m$ si existe un número $a^{-1}$ tal que
    \[
    a\cdot a^{-1} \equiv 1 \pmod m.
    \]
    \end{block}

    \vspace{0.25cm}
    \textbf{Interpretación}
    \begin{itemize}
        \item Multiplicar por $a^{-1}$ es ``deshacer'' el efecto de multiplicar por $a$ dentro del reloj.
        \item Si existe inverso, resolver $ax\equiv b$ es la versión de ``dividir por $a$'' dentro del reloj.
    \end{itemize}
\end{frame}

\begin{frame}{Por qué a veces no existe inverso}
    Compare dos números en reloj 12.

    \vspace{0.25cm}
    \textbf{Caso 1: $a=5$}\\
    Como vimos, $5\cdot 5 \equiv 1 \pmod{12}$, entonces 5 sí tiene inverso.

    \vspace{0.35cm}
    \textbf{Caso 2: $a=6$}\\
    Revisemos los múltiplos:
    \[
    6\cdot 1=6,\quad 6\cdot 2=12\rightarrow 0,\quad
    6\cdot 3=18\rightarrow 6,\quad 6\cdot 4=24\rightarrow 0.
    \]
    Solo aparecen posiciones 0 y 6. Nunca aparece 1.

    \vspace{0.3cm}
    \begin{alertblock}{Conclusión}
    Si nunca puedes llegar a 1, no hay inverso. En ese reloj, ``dividir por 6'' no existe.
    \end{alertblock}
\end{frame}

\begin{frame}{Regla con máximo común divisor}
    La regla que sí deben aprender es simple y poderosa.

    \vspace{0.35cm}
    \begin{block}{Criterio de inverso}
    $a$ tiene inverso en el reloj $m$ si y solo si
    \[
    \gcd(a,m)=1.
    \]
    \end{block}

    \vspace{0.25cm}
    \textbf{Cómo se usa}
    \begin{itemize}
        \item Si $\gcd(a,m)\neq 1$, multiplicar por $a$ colapsa posiciones y se pierde información.
        \item Si $\gcd(a,m)=1$, multiplicar por $a$ reordena posiciones sin colapsarlas, y existe el inverso.
    \end{itemize}
\end{frame}

\begin{frame}{Cuándo existe solución para $ax\equiv b$}
    Definimos $d=\gcd(a,m)$.

    \vspace{0.35cm}
    \begin{block}{Criterio de existencia}
    La ecuación $ax \equiv b \pmod m$ tiene solución si y solo si $d$ divide a $b$.
    \end{block}

    \vspace{0.25cm}
    \textbf{Lectura en español}
    \begin{itemize}
        \item Si $a$ y $m$ comparten un factor $d$, entonces $ax$ solo puede caer en ciertas posiciones.
        \item Solo se puede llegar a $b$ si $b$ es compatible con ese mismo factor.
    \end{itemize}
\end{frame}

\begin{frame}{Ejemplo clave en reloj 12}
    Resolvamos dos ecuaciones para ver el contraste.

    \vspace{0.3cm}
    \textbf{Caso A}
    \[
    6x \equiv 3 \pmod{12}.
    \]
    Aquí $\gcd(6,12)=6$, y 6 no divide a 3. No hay solución.

    \vspace{0.35cm}
    \textbf{Caso B}
    \[
    6x \equiv 6 \pmod{12}.
    \]
    Aquí 6 sí divide a 6. Hay soluciones.

    \vspace{0.25cm}
    Probando:
    \[
    x=1,3,5 \text{ funcionan, porque } (6x)\bmod 12 = 6.
    \]
\end{frame}

% -----------------------------------------------------------------------------
% SECCIÓN 5: PUENTE A SIMULACIÓN Y LCG
% -----------------------------------------------------------------------------
\section{Puente a la simulación}

\begin{frame}{De la aritmética de reloj al generador de uniformes}
    En el algoritmo de simulación, el primer paso fue:
    \[
    \text{generar } U\in[0,1).
    \]

    \vspace{0.25cm}
    \begin{block}{Lo que el computador genera primero}
    Genera estados enteros
    \[
    Z_n\in\{0,1,\dots,m-1\}
    \]
    y luego normaliza:
    \[
    U_n=\frac{Z_n}{m}.
    \]
    \end{block}

    \vspace{0.2cm}
    \begin{alertblock}{Conexión con auditoría}
    Si $Z_n$ cae en ciclos cortos o patrones, $U_n$ también.
    Entonces premios y tiempos heredarán esos patrones.
    \end{alertblock}
\end{frame}

\begin{frame}{Generador congruencial lineal}
    La idea es actualizar el estado con una regla de reloj.

    \vspace{0.25cm}
    \begin{block}{Regla general}
    Elegimos parámetros $a,c,m$ y una semilla $Z_0$:
    \[
    Z_{n+1} = (aZ_n + c)\bmod m,
    \qquad
    U_{n+1}=\frac{Z_{n+1}}{m}.
    \]
    \end{block}

    \vspace{0.2cm}
    \textbf{Lectura operacional}
    \begin{enumerate}
        \item multiplica el estado actual $Z_n$ por $a$,
        \item suma $c$,
        \item quédate con el residuo al dividir por $m$,
        \item divide por $m$ para obtener un número en $[0,1)$.
    \end{enumerate}
\end{frame}

\begin{frame}{Ejemplo malo: ciclo corto y valores repetidos}
    Usamos un reloj pequeño para poder auditar a mano:
    \[
    m=16,\quad a=5,\quad c=0,\quad Z_0=1.
    \]

    \vspace{0.15cm}
    \[
    \begin{array}{c|cccccccc}
    n & 0 & 1 & 2 & 3 & 4 & 5 & 6 & 7\\ \hline
    Z_n & 1 & 5 & 9 & 13 & 1 & 5 & 9 & 13\\
    U_n=\tfrac{Z_n}{16} & 0.0625 & 0.3125 & 0.5625 & 0.8125 & 0.0625 & 0.3125 & 0.5625 & 0.8125
    \end{array}
    \]

    \vspace{0.25cm}
    \begin{alertblock}{Qué pasa}
    El periodo es 4: la secuencia entra en un ciclo muy corto y repite los mismos cuatro valores.
    En un casino, esto puede volver predecibles premios y tiempos.
    \end{alertblock}
\end{frame}

\begin{frame}{Ejemplo bueno: periodo completo en un reloj pequeño}
    Ahora usamos parámetros que sí generan un recorrido largo:
    \[
    m=16,\quad a=5,\quad c=1,\quad Z_0=1.
    \]

    \vspace{0.15cm}
    \[
    \begin{array}{c|cccccccc}
    n & 0 & 1 & 2 & 3 & 4 & 5 & 6 & 7\\ \hline
    Z_n & 1 & 6 & 15 & 12 & 13 & 2 & 11 & 8\\
    U_n=\tfrac{Z_n}{16} & 0.0625 & 0.3750 & 0.9375 & 0.7500 & 0.8125 & 0.1250 & 0.6875 & 0.5000
    \end{array}
    \]

    \vspace{0.15cm}
    \[
    \begin{array}{c|cccccccc}
    n & 8 & 9 & 10 & 11 & 12 & 13 & 14 & 15\\ \hline
    Z_n & 9 & 14 & 7 & 4 & 5 & 10 & 3 & 0\\
    U_n=\tfrac{Z_n}{16} & 0.5625 & 0.8750 & 0.4375 & 0.2500 & 0.3125 & 0.6250 & 0.1875 & 0.0000
    \end{array}
    \]

    \vspace{0.2cm}
    \begin{exampleblock}{Lectura}
    En el paso siguiente vuelve a $Z_{16}=1$, así que el periodo es 16 (máximo para este reloj).
    \end{exampleblock}
\end{frame}

\begin{frame}{Condiciones para periodo completo}
    Para que el periodo sea máximo, existen condiciones clásicas.

    \vspace{0.25cm}
    \begin{block}{Teorema de Hull--Dobell}
    El periodo es $m$ si y solo si:
    \begin{enumerate}
        \item $\gcd(c,m)=1$,
        \item $a-1$ es divisible por todos los factores primos de $m$,
        \item si $4$ divide a $m$, entonces $4$ divide a $a-1$.
    \end{enumerate}
    \end{block}

    \vspace{0.2cm}
    \begin{exampleblock}{Chequeo rápido del ejemplo bueno}
        Para $m=16$, $a=5$, $c=1$:
        \[
        \gcd(1,16)=1,\quad a-1=4\text{ es divisible por }2,\quad 16\text{ es múltiplo de }4\text{ y }4\mid (a-1).
        \]
    \end{exampleblock}
\end{frame}

\begin{frame}{De $U$ a la simulación del casino}
    Ahora conectamos explícitamente con el algoritmo general.

    \vspace{0.25cm}
    \begin{block}{Paso 1: generar $U$}
    El generador produce $U_1,U_2,\dots$ en $[0,1)$.
    \end{block}

    \vspace{0.25cm}
    \begin{block}{Paso 2: transformar $U$ en variables del modelo}
    A partir de $U$ generamos:
    \begin{itemize}
        \item \textbf{premios discretos} por comparación con probabilidades acumuladas,
        \item \textbf{tiempos entre eventos} con transformada inversa,
        \item \textbf{eventos raros} con umbrales (por ejemplo $U>0.995$).
    \end{itemize}
    \end{block}

    \vspace{0.2cm}
    \textbf{Mensaje}\\
    Si el generador produce $U$ con patrones, esos patrones se heredan en premios y tiempos.
\end{frame}

\begin{frame}{Mini-auditoría: un premio discreto a partir de $U$}
    Supongamos que el tragamonedas define un premio $X$ con:
    \[
    \mathbb{P}(X=0)=0.90,\quad
    \mathbb{P}(X=10)=0.095,\quad
    \mathbb{P}(X=100)=0.005.
    \]

    \vspace{0.2cm}
    Acumuladas:
    \[
    F(0)=0.90,\quad F(10)=0.995,\quad F(100)=1.
    \]

    \vspace{0.25cm}
    \begin{block}{Regla de generación}
    \begin{itemize}
        \item Si $U\le 0.90$, devuelve $X=0$.
        \item Si $0.90<U\le 0.995$, devuelve $X=10$.
        \item Si $U>0.995$, devuelve $X=100$.
    \end{itemize}
    \end{block}

    \vspace{0.15cm}
    \textbf{Lectura para auditoría}\\
    Si el generador subrepresenta valores altos de $U$, el premio grande aparecerá menos de lo anunciado.
\end{frame}

% -----------------------------------------------------------------------------
% SECCIÓN 6: TRANSFORMADA INVERSA
% -----------------------------------------------------------------------------
\section{Transformada inversa}

\begin{frame}{Del uniforme a las variables del modelo}
    El generador entrega $U$ en $[0,1)$.
    Pero el casino necesita variables con distribuciones específicas.

    \vspace{0.25cm}
    \begin{itemize}
        \item tiempos entre eventos,
        \item premios discretos,
        \item montos con colas pesadas,
        \item probabilidades de eventos raros.
    \end{itemize}

    \vspace{0.35cm}
    \begin{block}{Objetivo}
    Construir una regla que, a partir de $U$, genere una variable $X$ con la distribución deseada.
    \end{block}
\end{frame}

\begin{frame}{La CDF como mapa de probabilidad acumulada}
    La función de distribución acumulada es:
    \[
    F(x)=\mathbb{P}(X\le x).
    \]

    \vspace{0.25cm}
    \textbf{Interpretación}
    \begin{itemize}
        \item $F(x)$ es la fracción de probabilidad que queda a la izquierda de $x$.
        \item Si $U$ es una fracción entre 0 y 1, buscamos el punto $x$ donde se acumula esa fracción.
    \end{itemize}

    \vspace{0.2cm}
    \begin{center}
    \begin{tikzpicture}[scale=0.95]
        \draw[-{Stealth[length=3mm]}] (0,0) -- (6.2,0) node[right] {$x$};
        \draw[-{Stealth[length=3mm]}] (0,0) -- (0,3.2) node[above] {$F(x)$};
        \draw[thick] (0,0) .. controls (2,0.2) and (3,2.2) .. (6,3);
        \draw[dashed] (0,2.1) -- (3.7,2.1);
        \draw[dashed] (3.7,0) -- (3.7,2.1);
        \node at (0.55,2.25) {$U$};
        \node at (4.2,-0.25) {$x$};
    \end{tikzpicture}
    \end{center}
\end{frame}

\begin{frame}{Regla de la transformada inversa}
    Si $U\sim \text{Uniforme}(0,1)$ y queremos que $X$ tenga CDF $F$:

    \vspace{0.35cm}
    \begin{block}{Definición}
    \[
    X = F^{-1}(U),
    \qquad
    F^{-1}(u)=\inf\{x:\ F(x)\ge u\}.
    \]
    \end{block}

    \vspace{0.25cm}
    \textbf{Lectura en español}
    \begin{itemize}
        \item $U$ escoge un percentil.
        \item $F^{-1}$ devuelve el valor de $x$ correspondiente a ese percentil.
    \end{itemize}
\end{frame}

\begin{frame}{Por qué funciona}
    Queremos que la CDF de $X$ sea exactamente $F$.

    \vspace{0.25cm}
    Tomemos cualquier $x$:
    \[
    \mathbb{P}(X\le x)
    =
    \mathbb{P}(F^{-1}(U)\le x).
    \]

    \vspace{0.15cm}
    Por la definición de inversa generalizada:
    \[
    F^{-1}(u)\le x \ \Longleftrightarrow\ u\le F(x).
    \]

    \vspace{0.15cm}
    Entonces:
    \[
    \mathbb{P}(X\le x)=\mathbb{P}(U\le F(x))=F(x),
    \]
    porque $U$ es uniforme.
\end{frame}

\begin{frame}{Transformada inversa en el caso discreto}
    Premios discretos: valores $\{x_1<x_2<\cdots<x_k\}$ con probabilidades $p_1,\dots,p_k$.

    \vspace{0.25cm}
    Definimos acumuladas:
    \[
    F(x_j)=\sum_{i=1}^j p_i.
    \]

    \vspace{0.25cm}
    \begin{block}{Algoritmo}
    Genera $U$ y toma el menor $j$ tal que $F(x_j)\ge U$.
    Devuelve $X=x_j$.
    \end{block}

    \vspace{0.25cm}
    \begin{exampleblock}{Ejemplo}
    $x=\{0,10,100\}$ con $p=\{0.90,0.095,0.005\}$.
    \begin{itemize}
        \item Si $U\le 0.90$, devuelve 0.
        \item Si $0.90<U\le 0.995$, devuelve 10.
        \item Si $U>0.995$, devuelve 100.
    \end{itemize}
    \end{exampleblock}
\end{frame}

\begin{frame}{Ejemplo continuo: exponencial}
    Para modelar tiempos entre eventos, una elección frecuente es la exponencial con tasa $\lambda$:
    \[
    F(x)=1-e^{-\lambda x},\qquad x\ge 0.
    \]

    \vspace{0.25cm}
    Igualamos $U=F(x)$ y despejamos:
    \[
    U=1-e^{-\lambda x}
    \Rightarrow
    e^{-\lambda x}=1-U
    \Rightarrow
    x=-\frac{1}{\lambda}\ln(1-U).
    \]

    \vspace{0.25cm}
    \begin{block}{Fórmula práctica}
    \[
    X=-\frac{1}{\lambda}\ln(U).
    \]
    \end{block}
\end{frame}

% -----------------------------------------------------------------------------
% SECCIÓN 7: CIERRE
% -----------------------------------------------------------------------------
\section{Cierre}

\begin{frame}{El flujo completo del motor}
    \begin{block}{Cadena de construcción}
    \[
    Z_{n+1}=(aZ_n+c)\bmod m
    \ \Rightarrow\
    U_n=\frac{Z_n}{m}
    \ \Rightarrow\
    X=F^{-1}(U_n).
    \]
    \end{block}

    \vspace{0.25cm}
    \textbf{Lectura para auditoría}
    \begin{itemize}
        \item La parte del reloj decide si $U$ es poco predecible y estadísticamente útil.
        \item La transformada decide si el comportamiento del juego coincide con las probabilidades declaradas.
    \end{itemize}
\end{frame}

\begin{frame}{Ejercicios}
    \begin{enumerate}
        \item En reloj 10, liste los residuos de $4x$ para $x=0,1,\dots,9$. ¿Puede aparecer el residuo 1?
        \item En reloj 12, encuentre un $x$ tal que $5x$ deje residuo 7.
        \item En el ejemplo bueno con $m=16,a=5,c=1$, calcule $Z_1,\dots,Z_5$ y sus $U_n$.
        \item Premios $\{0,10,100\}$ con probabilidades $\{0.90,0.095,0.005\}$.\\
              Si $U=0.25$, ¿qué premio sale? Si $U=0.998$, ¿qué premio sale?
    \end{enumerate}
\end{frame}

\begin{frame}{Referencias}
    \small
    \begin{itemize}
        \item Banks, J., Carson, J. S., Nelson, B. L., \& Nicol, D. M. \textit{Discrete-Event System Simulation}.
        \item Law, A. M. \textit{Simulation Modeling and Analysis}.
    \end{itemize}
\end{frame}

\end{document}
